\section{Обзор источников}

Рассматриваемый в работе набор данных iPinYou получил широкое распространение в академических исследованиях по вычислительной рекламе и машинному обучению, поскольку представляет собой редкий пример публично доступных логов RTB-аукционов, включающих события на разных этапах рекламной воронки (запрос ставки, факт показа, клик и конверсия). В частности, данные iPinYou активно используются для задач предсказания отклика пользователя (CTR/CVR), оптимизации назначения ставок на стороне DSP, а также для сравнения стратегий закупки рекламного инвентаря в условиях ограниченного бюджета~\cite{rtb-research, rtb-benchmarking}.

В обзорной работе~\cite{rtb-research} авторы рассматривают RTB как отдельное направление исследований в медийной рекламе и систематизируют ключевые задачи, возникающие на стороне DSP. Существенный акцент делается на том, что решение задачи назначения ставки невозможно свести к одной модели: оно включает (1) оценку вероятности отклика пользователя (например, предсказание клика/конверсии), (2) принятие решения о ставке с учётом бюджета и цели кампании, (3) учёт ограничений и особенностей аукционного механизма. Отдельно отмечаются практические сложности, характерные для RTB: сильная разреженность целевых событий (клики и особенно конверсии), необходимость работы в строгих временных ограничениях. Таким образом, \cite{rtb-research} задаёт концепт, в котором связаны предсказательные модели машинного обучения (CTR/CVR) и оптимизационные задачи RTB (назначение ставок, распределение бюджета).

Работа \cite{rtb-benchmarking} посвящена построению протокола сравнения алгоритмов назначения ставок на основе набора данных iPinYou и фактически является одной из отправных точек для последующих исследований. Авторы проводят первичный анализ структуры и свойств датасета, формулируют постановку задачи DSP Bid Optimization Problem и предлагают подход к офлайн-оценке стратегий на логах аукционов (симуляция закупки при фиксированном бюджете). Далее рассматривается 2 этапа: сначала обучается модель предсказания вероятности клика и/или конверсии (pCTR/pCVR), затем на основе полученной оценки ценности показа применяется стратегия назначения ставки. В экспериментальной части сравниваются базовые стратегии (в частности, константная ставка и случайное назначение ставок) и более содержательные правила, использующие предсказанную вероятность отклика. Показано, что использование моделей pCTR/pCVR и корректно выбранной стратегии ставок позволяет существенно повысить итоговые KPI по сравнению с наивными базовыми подходами.
